\section{How long would a solve take?}
\label{sec:solvetime}

The rates in this paper translate directly into calendar time and
money, and this section does that translation for every device we
measured. Nothing here is a new measurement: each estimate divides the
expected work by a rate recorded earlier in the paper and states what
else it assumes. The question for ECC2K-130 is no longer whether a solve
is feasible but how many cards, for how long, at what price. For the
NIST curves, the same arithmetic gives an answer that no plausible
amount of hardware changes.

\subsection{The model}
\label{sec:solvetime-model}

Three quantities set the time to a collision.

\paragraph{The work.} The $\sigma$-walk needs $2^{60.9}\approx2.15\cdot10^{18}$
iterations in expectation (Section~\ref{sec:walk-constant}). One
device running at $R$ billion iterations per second therefore needs
$2.15\cdot10^{18}/R$ nanoseconds, or $68.2/R$ years
(Table~\ref{tab:history-compare}), and a fleet of $N$ such devices
needs $1/N$ of that.

\paragraph{The spread.} Rho's running time is random. The probability
of no collision after a fraction $c$ of the expected work is about
$\exp(-\pi c^2/4)$ (Section~\ref{sec:campaign-state}). Half of all
solves finish by $c=0.94$, nine in ten by $c=1.71$ and 99 in 100 by
$c=2.42$ (our arithmetic). So a plan that must finish with $90\%$
probability should budget $1.7$ times the expected work.

\paragraph{The detection delay.} A collision is only seen when the
second of the two merged trails reaches a distinguished point. At
cutoff 32 that is on average another $2^{28.41}$ steps of one walk
(Section~\ref{sec:campaign-interval}). Each walk advances at $R$
divided by the number of walks on the card. With $6.16$ million walks
on an RTX~PRO~6000, the delay is about $42$ hours at $14.64$ and $38$
hours at $15.88$, whatever the fleet size (our arithmetic; the campaign
note gives $43$ GPU-hours per GPU at $14.11$~\evcite{\evref{ops}}).
Every card keeps walking during the delay, so a large fleet pays for it
in GPU-hours as well as calendar time. At 1{,}000 cards it adds about
as much work as the solve itself. Cutoff 34 cuts the delay to about
$4.4$ hours at $15.88$, at the cost of an $8.8\times$ larger corpus
(about $1.7$~TB instead of $0.19$~TB).

A fleet of $N$ devices at $R$ therefore expects a collision after
\[
  T(N) \approx \frac{2^{60.9}}{N\,R} + D,
\]
where $D$ is the detection delay. The first term is multiplied by $0.94$
for the median and by $1.71$ for a $90\%$ plan.

\paragraph{The corpus.} The campaign has already walked $2^{56.68}$
iterations without a collision (Table~\ref{tab:state}). A new fleet
that runs the same $\sigma$-walk at cutoff 32 and merges with that
corpus inherits the credit. Given no collision so far, the expected
remaining work is $0.949$ of the total (our arithmetic). That is a
saving of $5\%$; it changes the plan, not the scale.

\subsection{One device, a hundred, a thousand}
\label{sec:solvetime-table}

Table~\ref{tab:solvetime} applies the first term of the model to every
rate in this paper. The $\sigma$-walk rows are the survey and paired
rates of Table~\ref{tab:survey}. Points from these rows can merge with
the campaign corpus. The other rows use other walks or other platforms.

\begin{table}[t]
\centering
\footnotesize
\setlength{\tabcolsep}{4pt}
\caption{Expected time to walk $2^{60.9}$ iterations of ECC2K-130 at
  each recorded rate, our arithmetic. ``Years, one'': device-years on a
  single device. ``Days, 100'' and ``Days, 1{,}000'': calendar days for
  a fleet of that size, excluding the detection delay $D$ (about
  $1.6$--$1.7$ days on the RTX~PRO~6000 at cutoff 32). The median is
  $0.94\times$ these figures and the $90\%$ point $1.71\times$. ``Cost'':
  the expected work at the Modal list price implied by
  Table~\ref{tab:survey}'s iterations per dollar (2026-09-19). Spot
  prices are much lower (text). Table-walk rates are absent because no
  table-walk build has a measured cost per solve
  (Section~\ref{sec:sigma-vs-table}).}
\label{tab:solvetime}
\begin{tabular}{@{}llrrrrr@{}}
\toprule
device & build & billion it/s & years, one & days, 100 & days, 1{,}000 & cost \\
\midrule
\multicolumn{7}{@{}l}{\emph{$\sigma$-walk (merges with the campaign corpus)}} \\
RTX PRO 6000 & campaign preset & 14.64 & 4.66 & 17.0 & 1.70 & \$124k \\
RTX PRO 6000 & nibble-table $\sigma$ & 15.88 & 4.29 & 15.7 & 1.57 & \$114k \\
L40S & survey & 8.70 & 7.84 & 28.6 & 2.86 & \$134k \\
B200 & survey & 8.84 & 7.71 & 28.2 & 2.82 & \$423k \\
B300 & survey & 7.58 & 9.00 & 32.9 & 3.29 & \$560k \\
H200 & survey & 7.63 & 8.93 & 32.6 & 3.26 & \$356k \\
H100 & survey & 7.54 & 9.04 & 33.0 & 3.30 & \$313k \\
A100-80GB & survey & 4.63 & 14.7 & 53.8 & 5.38 & \$323k \\
L4 & survey & 2.56 & 26.6 & 97.2 & 9.72 & \$186k \\
A10 & survey & 2.42 & 28.1 & 102.8 & 10.3 & \$272k \\
T4 & no \texttt{clmad} & 0.541 & 126 & 460 & 46.0 & \$652k \\
\addlinespace
\multicolumn{7}{@{}l}{\emph{other walks and platforms}} \\
RTX PRO 4500 & $x$-only bridge map & 12.50 & 5.46 & 19.9 & 1.99 & --- \\
RTX PRO 4500 & $\sigma$-walk & 5.11 & 13.3 & 48.8 & 4.88 & --- \\
VU47P (AWS F2) & 128 engines & 8.25 & 8.26 & 30.2 & 3.02 & --- \\
Xeon, one core & AVX-512 bitsliced & 0.0126 & 5{,}410 & --- & --- & --- \\
\addlinespace
\multicolumn{7}{@{}l}{\emph{bound and history}} \\
RTX PRO 6000 & ceiling (Section~\ref{sec:bounds-6000}) & 24--26 & 2.6--2.8 & 9.6--10.4 & 0.96--1.04 & --- \\
GTX 295 card & 2010 kernel & 0.064 & 1{,}065 & 3{,}891 & 389 & --- \\
\bottomrule
\end{tabular}
\end{table}

% Expected days = 2^60.9 / (N R) / 86400 = 24896 / (N R) for R in billion it/s.
% Plotted parametrically in log10(N) so the curves sample evenly on the log axis.
\begin{figure}[t]
\centering
\begin{tikzpicture}
\begin{loglogaxis}[
  width=0.92\textwidth, height=0.56\textwidth,
  xmin=1, xmax=10000, ymin=0.3, ymax=8e5,
  xlabel={devices in the fleet},
  ylabel={expected days to a collision},
  xmajorgrids, ymajorgrids,
  legend style={at={(0.5,-0.2)},anchor=north,font=\scriptsize},
  legend columns=3, legend cell align=left,
  every axis plot/.append style={thick},
]
\addplot[cA,domain=0:4,samples=2] ({10^x},{24896/(15.884*10^x)});
\addlegendentry{RTX PRO 6000, $\sigma$ 15.88}
\addplot[cA,dashed,domain=0:4,samples=60] ({10^x},{24896/(15.884*10^x)+1.6});
\addlegendentry{same, with the cutoff-32 delay}
\addplot[cB,densely dashed,domain=0:4,samples=2] ({10^x},{24896/(25*10^x)});
\addlegendentry{RTX PRO 6000 ceiling, 25}
\addplot[cE,domain=0:4,samples=2] ({10^x},{24896/(8.840*10^x)});
\addlegendentry{B200, $\sigma$ 8.84}
\addplot[cD,domain=0:4,samples=2] ({10^x},{24896/(7.540*10^x)});
\addlegendentry{H100, $\sigma$ 7.54}
\addplot[cC,domain=0:4,samples=2] ({10^x},{24896/(2.561*10^x)});
\addlegendentry{L4, $\sigma$ 2.56}
\addplot[inkS!60,domain=0:4,samples=2] ({10^x},{24896/(0.541*10^x)});
\addlegendentry{T4, no \texttt{clmad}, 0.541}
\addplot[inkS,densely dotted,domain=0:4,samples=2] ({10^x},{24896/(0.064*10^x)});
\addlegendentry{GTX 295 card, 2010, 0.064}
% calendar references
\draw[gridc!60!inkS,line width=0.4pt] (axis cs:1,7) -- (axis cs:10000,7);
\node[anchor=north west,font=\tiny,text=inkS,inner sep=1.5pt] at (axis cs:1,7) {one week};
\draw[gridc!60!inkS,line width=0.4pt] (axis cs:1,30) -- (axis cs:10000,30);
\node[anchor=north west,font=\tiny,text=inkS,inner sep=1.5pt] at (axis cs:1,30) {one month};
\draw[gridc!60!inkS,line width=0.4pt] (axis cs:1,365) -- (axis cs:10000,365);
\node[anchor=north west,font=\tiny,text=inkS,inner sep=1.5pt] at (axis cs:1,365) {one year};
\end{loglogaxis}
\end{tikzpicture}
\caption{Expected calendar time to walk $2^{60.9}$ iterations of
  ECC2K-130 against fleet size, at measured one-device rates in billions
  of iterations per second (Table~\ref{tab:solvetime}), with the 2010
  GTX~295 for comparison. Solid lines count walking time only. The
  dashed line adds the RTX~PRO~6000's detection delay at cutoff 32
  (about $1.6$ days), which sets a floor that more cards cannot lower.
  The median solve is $0.94\times$ the plotted time and the $90\%$
  point $1.71\times$.}
\label{fig:solvetime}
\end{figure}


On the best card, a solve is about four device-years. A fleet of five
RTX~PRO~6000s finishes in about a year. With 56 the expected time is a
month, including the detection delay. With 290 it is a week. A thousand
cards bring the expected time down to about $3.2$ days, $1.6$ of
walking and $1.6$ of waiting for the last trail to reach a
distinguished point. Nine times in ten, a thousand cards finish within
$4.3$ days. Beyond a few hundred cards the delay dominates, which is
where Figure~\ref{fig:solvetime}'s dashed line bends away from the
solid one. A larger fleet buys little calendar time unless it also
lowers the cutoff.

The new GPUs moved this by more than two orders of magnitude. With the
2010 kernel on a GTX~295, a one-year solve needed about 1{,}066 cards
and a one-week solve more than $55{,}000$. Today the same one-year and
one-week targets need five cards and 290. If the RTX~PRO~6000 reached
its ceiling, the solve would be about $2.7$ device-years. Even the
campaign's own peak, about $83$~billion iterations/s from a mixed spot
fleet (Table~\ref{tab:state}), would finish the remaining expected
work in about $285$ days.

\paragraph{Cost.} The cost of the work does not depend on the size of
the fleet. At Modal list prices the cheapest per solve are the
RTX~PRO~6000 and the L40S, about \$114k--\$134k. The data-centre
Blackwell and Hopper parts cost two and a half to five times as much
for the same work, because this integer-bound walk does not use what
they are priced for (Section~\ref{sec:survey}). Spot capacity is
cheaper again. At the 2026-09-11 spot price of \$1.31 per RTX~PRO~6000
hour the expected work is about \$56k at the campaign's collection rate
of $14.11$. An AWS F2 at its spot price comes to about \$51k
(Section~\ref{sec:fpga-vs-gpu}). The best quote seen, an L4 in
\texttt{eu-north-1} on 2026-09-18, priced a full walk at about
\$38k~\evcite{\evref{ops}; \evref{quotes}}. Those spot figures hold only
while spot capacity lasts. The spot market reclaimed the campaign's
whole fleet on 2026-09-18 (Section~\ref{sec:campaign-state}), and a
solve on spot capacity has to plan for that.

\paragraph{What these estimates assume.} Every rate in
Table~\ref{tab:solvetime} is sustained for the whole solve. The survey
rows are single unpaired runs, and the campaign's own fleet measured
its RTX~PRO~6000 workers at a median of $14.66$
(Table~\ref{tab:cutoff}), so the survey rates are close to what a fleet
sees but are not guaranteed. The bridge map on the RTX~PRO~4500 walks a
different function, so its points do not merge with the
$\sigma$-walk corpus. Its frozen study measured $1.007\times$ the work
of the map it replaced (Section~\ref{sec:walk-bridge}). The FPGA row is
one card of an \texttt{f2.6xlarge}. The CPU row is one core and assumes
the per-core rate holds across a server. The costs exclude storage,
egress and the merge.

\subsection{The NIST curves on the same footing}
\label{sec:solvetime-nist}

The same model applies to the NIST walk of Section~\ref{sec:nist}, and
it shows why that section's rates say nothing about the security of
those curves. With the negation map, the expected work for a group of
prime order $n$ is about $\sqrt{\pi n/4}$ iterations. For P-256 that is
$2^{127.8}\approx3.0\cdot10^{38}$ iterations, and for P-384
$2^{191.8}\approx5.6\cdot10^{57}$ (our arithmetic from the standard group
orders; any walk constant above one only adds to it).
\begin{itemize}[nosep]
\item \textbf{P-256.} At the best measured $1.632$~billion
  iterations/s on one RTX~PRO~6000, the expected time is $1.9\cdot10^{29}$
  seconds, about $5.9\cdot10^{21}$ card-years. A billion cards would
  need $5.9\cdot10^{12}$ years, more than 400 times the age of the
  universe.
  At the multiplies-only bound of $12.2$ (Section~\ref{sec:bounds-nist})
  it is still $7.8\cdot10^{20}$ card-years.
\item \textbf{P-384.} At $0.548$ the expected time is $1.0\cdot10^{49}$
  seconds, about $3.2\cdot10^{41}$ card-years.
\end{itemize}
The engineering of Section~\ref{sec:nist} multiplied the P-256 rate by
$1.52$ and the P-384 rate by $2.61$, and divided these times by the
same factors. P-256 needs about $2^{67}$ times the work of ECC2K-130, so
a speedup of any engineering size leaves it where it was. The levers
that bring ECC2K-130 within a fleet's reach shorten a P-256 attack by a
factor that disappears against the exponent.
